\section{Reproducibility and model terms}
\label{c-app:reproducibility}

The main panel contains seven instruction-tuned checkpoints, with CodeLlama-34B as a
scale check. We use CodeLlama \citep{roziere2023codellama}, CodeGemma
\citep{codegemmateam2024codegemma}, Llama~3 \citep{grattafiori2024llama3}, Gemma~3
\citep{gemmateam2025gemma3} and OLMo~2 \citep{olmo2024olmo2,walsh2025olmo2}.
The model identifiers are \texttt{codellama/CodeLlama-7b-Instruct-hf} and its 34b variant,
\texttt{google/codegemma-7b-it}, \texttt{meta-llama/Llama-3.2-3B-Instruct},
\texttt{meta-llama/Llama-3.1-8B-Instruct}, \texttt{google/gemma-3-4b-it},
\texttt{allenai/OLMo-2-0425-1B-Instruct} and \texttt{allenai/OLMo-2-1124-7B-Instruct}.
These use the Llama~2/CodeLlama, Llama~3.1 or 3.2 Community Licenses, Gemma Terms of Use,
and Apache~2.0, respectively. The study uses their released weights for research inference
and does not redistribute weights. CRUXEval uses the MIT license.

Synthetic data consist of digits, generated identifiers and short Python functions,
without personal records or human annotations. Each tested level has 1,000 matched
problems per target. Main behavior uses three sets of three demonstrations with seeds
7, 11 and 13; programs and demonstration identities remain fixed across compared formats.
Greedy decoding uses the full vocabulary. Parsers retain failed outputs in denominators;
trace-error measurements use the first assignment, while final-answer accuracy scores
the reported return value. Probe training and testing are disjoint in arithmetic computation,
ignoring identifier spelling. The independent-name replication uses separate screened
names and 200 new computations, as described in Appendix~\ref{c-app:round6-identifiers}.

Probes use full-batch SGD with learning rate $10^{-3}$ for 10,000 epochs and seeds 0, 1 and 2.
Five-fold layer selection is disjoint in computation; bootstrap intervals condition on
those selectors. These are linear-predictor fits, not language-model training.
Combined recorded allocations across the arithmetic and code experiments total
\NUM{c-compute-shared-gpuh} GPU-hours on NVIDIA A30, H100 and H200 GPUs, an upper bound
including idle and shared workers. Software is PyTorch~2.11.0 (CUDA~12.9)
\citep{paszke2019pytorch}, Transformers~5.14.1 \citep{wolf2020transformers},
NumPy~2.3.5 and scikit-learn~1.9.1 \citep{pedregosa2011sklearn}.
